\section{引言}
\label{sec:intro}

多模态大语言模型（MLLM）已推动视觉问答、图像描述与开放式多模态推理快速发展，但即使是强模型仍易出现幻觉：生成流畅却接地不足、虚构图像中不存在的物体，或在强语言先验下对视觉不支持的实体过度自信。当正确证据稀疏、局部化或被无关视觉上下文淹没时，这类失败尤为有害。

近期工作表明，\textbf{推理期干预}是替代重训宿主模型的有效途径。对比解码（如 VCD）、基于注意力的控制（如 OPERA）以及层间对比解码（如 DoLa）均将自回归解码视为可操作的干预面。然而，仍存在更严格的问题：

\begin{quote}
\textbf{解码期缓解能否在异构 MLLM 家族间保持可迁移性——无需按模型单独训练，且不依赖宿主内部的隐状态、logit 或注意力？}
\end{quote}

该目标不同于单点提升骨干性能。任何消费模型特异多模态激活、并在宿主原生几何中写入校正的方法，即便视觉—语言权重冻结，仍属\textbf{宿主耦合}。\textbf{隐状态校准}类路线说明了这一现象，但难以满足严格的即插即用可迁移要求。

因此我们提出\textbf{可迁移性边界}：机制若完全运行于模型原生的隐空间、logit 空间或注意力空间内，则无法实现严格意义上的热插拔。可检验的替代方案是将所有评分限制在由\textbf{冻结的公开视觉特征}、\textbf{冻结文本编码器}、\textbf{解码得到的候选片段}以及\textbf{确定性}分词器桥接所构成的\textbf{共享外部空间}中，且\textbf{不}引入绑定于特定宿主的可学习适配器。

\textbf{CHORD}（Calibrating Hallucinations via Object-Resonant Decoding）体现该原则。它执行\textbf{候选驱动的视觉共振}：关键信号并非未完成前缀与图像的全局相似度，而是\textbf{当前候选片段}是否激活一致的局部视觉证据，以及该证据是否支持干预。缓解被刻画为对与候选绑定的外部、物体级证据的选择性调用，而非对无关区域或纯语言延续偏置的整体重分配。

由此有两点推论：其一，除原始幻觉降幅外，还应在\textbf{外部编码器与区域库固定、仅分词器桥接随宿主确定性变化}的前提下，考察 CHORD 在\textbf{不同宿主家族}上是否\textbf{稳定}；其二，基于稠密区域库的片段级共振应明显优于朴素的全局图像先验，否则额外机制难以自证必要。

本文贡献如下：
\begin{enumerate}
  \item 形式化限制宿主内部解码器严格可移植控制的\textbf{可迁移性边界}。
  \item 提出\textbf{CHORD}：\textbf{免训练}解码框架，在外部嵌入空间中利用物体共振证据与闭式\textbf{共振增量 $\Delta S$}校准活跃 logit 前沿。
  \item 引入\textbf{无参数分词器桥接}及\textbf{连续 Top-$1$ 监控}（配合全 $K$ 回退），使共振检验轻量，并在不依赖朴素熵门控的前提下抑制过度自信幻觉。
  \item 在 POPE、CHAIR 上给出含解码基线、消融与效率分析的\textbf{评估协议}。
\end{enumerate}

全文结构：第~\ref{sec:related} 节为相关工作；第~\ref{sec:method} 节给出边界与 CHORD 方法；第~\ref{sec:exp} 节为实验；第~\ref{sec:conclusion} 节为结论。
